\section{Prices for the inherited activities}
\label{sec:inherited-activity-prices}
The remaining logarithmic activities come from the incoming
representation: its regular activities, cubic and quartic path
terms, quadratic loop terms, and the true long-path tails on the
fine and coarse levels. We convert their existing coefficient
bounds into prices on the full coarse read supports. This section
assumes those incoming coefficient bounds on the same representation
and paths as the selected step. It proves the price conversion;
maintenance of the incoming bounds along the renormalisation run
is a separate construction requirement.

\subsection{Coarsening a connected cover}
Let $\pi$ be the parent map with block length $L$, and let
$c_M(X)$ be the connected covering number at radius $M$. Thus
$X$ is contained in the union of $M$-balls around $c_M(X)$
centres, and that union is connected by unit lattice steps.
For $B\ge1$, $M\le2B$ and integer $\ell\ge0$, put
\[
 P_\ell=b_{\ell+6},\qquad v_\ell=10P_\ell/L,
 \qquad Y=\operatorname{halo}_{\ell B}(\pi X).
\]
Then
\begin{equation}\label{eq:inherited-cover-conversion}
 c_B(Y)\le P_\ell+v_\ell c_M(X).
\end{equation}

Here is a proof retaining the factor $L^{-1}$. Choose a minimal
cover with centre set $D$. The graph joining centres at distance
at most $2M+1$ is connected. Traverse a spanning tree out and
back, replacing each edge by steps of max-distance at most one
(diagonal steps are allowed). This gives a walk
visiting all centres, of length at most $(4M+2)|D|$.
Sample it every $LB$ steps and apply $\pi$. If its length is $N$,
the resulting centre set $D'$ satisfies
$LB|D'|\le N+LB$. The parent-distance inequality
\[
 L\,d(\pi x,\pi y)\le d(x,y)+L
\]
shows that every projected walk point is within $2B$ of a
sampled centre and that consecutive sampled centres are within
$2B$. Every point of $X$ is within $M\le2B$ of an original
centre; its parent is therefore within $3B$ of that centre's
parent. Thus $Y$ is covered by the $(\ell+5)B$-balls about
$D'$, with connected union. Tile their slightly enlarged
$(\ell+6)B$-balls by $B$-balls with centres on a $B$-spaced
four-dimensional grid. There are at most $P_\ell$ tiles per
centre. These enlarged halos overlap along the sampled walk,
so the resulting cover remains connected. Finally
$(4M+2)|D|\le10B|D|$ proves
\eqref{eq:inherited-cover-conversion}. For an empty set the
inequality is immediate. Wrapping any of these walks or grids
only decreases the number of distinct centres.

If $o\in X$ and $y\in Y$, the same spanning walk gives
\begin{equation}\label{eq:inherited-anchor-distance}
 d(y,\pi o)\le(\ell+20)B\{1+c_M(X)\}.
\end{equation}
Indeed two points of $X$ are at distance at most
$(4M+2)c_M(X)+2M$; apply the parent-distance inequality and
then add the $\ell B$ halo. The displayed larger constant
bounds the result for $M\le2B$, $B\ge1$.
Inspecting the fine plaquettes touching a support adds at most
four fine lattice steps. Their parents lie within five coarse
steps, so replacing $\ell$ by $\ell+5$ pays all these reads.

\subsection{From an anchor bound to a site bound}
We give a general counting estimate. Suppose nonnegative weights
$b_i$, integer costs $d_i\ge0$, fine anchors $o_i$, and coarse
supports $Y_i$ satisfy
\[
 \sup_o\sum_{i:o_i=o}b_i e^{\sigma d_i}\le\Gamma,
 \qquad c_B(Y_i)\le v+v'd_i,
 \qquad y\in Y_i\Longrightarrow
 d(y,\pi o_i)\le W(1+d_i).
\]
Let $k,W$ be nonnegative integers and
$0\le g_i\le C b_i(1+d_i)^k$ with $C\ge0$, and assume
\begin{equation}\label{eq:inherited-reserve}
 a\ge0,\quad\varepsilon>0,\quad av'+2\varepsilon\le\sigma.
\end{equation}
For a fixed coarse site $y$, the fine anchors of labels with
$d_i=n$ and $y\in Y_i$ have parents in a radius-$W(1+n)$
ball. There are at most $L^4b_W(1+n)^4$ such anchors, because
a parent has exactly $L^4$ residues. The sum of $g_i e^{a c_B(Y_i)}$
over these labels is at most
\[
 L^4b_W\Gamma C e^{av}(1+n)^{4+k}e^{-2\varepsilon n}.
\]
Use the exponential series at $\varepsilon(1+n)$ to obtain
\[
 (1+n)^{4+k}\le
 \varepsilon^{-(4+k)}(4+k)!e^\varepsilon e^{\varepsilon n}.
\]
Summing the geometric series proves
\begin{equation}\label{eq:inherited-payment}
 \sup_y\sum_{i:y\in Y_i}g_i e^{a c_B(Y_i)}
 \le e^{av}\Gamma C L^4b_W
 \frac{\varepsilon^{-(4+k)}(4+k)!e^\varepsilon}
      {1-e^{-\varepsilon}}.
\end{equation}
Restricting to $d_i\ge n_0$ adds the factor
$e^{-\varepsilon n_0}$. The proof works with arbitrary label
multiplicities, since the input already sums all labels at an
anchor. Nonnegative finite partial sums justify the estimate
also for a countable alphabet.

We need an additional factor $e^{2\varepsilon d_i}$ in the
output price. Apply the preceding argument to weights
$b_i e^{2\varepsilon d_i}$ and exponent $\sigma-2\varepsilon$.
The sufficient condition becomes $av'+4\varepsilon\le\sigma$.
For inherited supports use
\eqref{eq:inherited-cover-conversion} with
$d_i=c_M(X_i)$ and $W=(\ell+20)B$.
If $4\varepsilon<\sigma$, choose $L$ large enough that
$a10P_\ell/L+4\varepsilon\le\sigma$, before increasing $x$.
The incoming cover radius $M=B(h)$ is eventually at most $2B(x)$:
$h=x+c_L$ with $c_L$ fixed, and the logarithmic-square radius
has $B(x+c_L)/B(x)\to1$. All factors in
\eqref{eq:inherited-payment} other than $\Gamma C$ then have
upper exponent zero.

\subsection{Path lengths, including winding paths}
For a bundle of $r$ paths based at one fine anchor, let $d_i$ be
the sum of their true lengths, including their final plaquette
walks. Dividing each walk into pieces of $B$ steps shows that its
visited sites have connected $B$-cover number at most
$1+\text{length}/B$. All paths share the anchor, so their union
with the anchor has cover at most $1+r+d_i/B$. Applying
\eqref{eq:inherited-cover-conversion} at $M=B$ gives
\[
 c_B(Y_i)\le P_\ell+v_\ell(1+r)+(v_\ell/B)d_i.
\]
The support lies within $(\ell B+1)(1+d_i)$ of the parent
anchor. Thus the preceding site estimate applies with
$v'=v_\ell/B$. For any fixed $L$, this coefficient tends to
zero, so the required strict exponential reserve is eventually
available. No path length is replaced by its endpoint distance;
in particular, winding and repeated edges retain their costs.

Suppose the path coefficients satisfy
$\sum_i b_i e^{\sigma d_i}\le\Gamma$ at each translated
anchor. An order-$r$ atom evaluated with its original small-field
mask has absolute value at most $b_i\tau^r$: every one of its
$r$ curvature arguments has norm at most $\tau$, while the
unhealthy branch is zero. An external multiplier bounded by
$D(x)$ therefore has scalar price $D(x)\tau^r b_i$.
The site conversion preserves the sum of the upper exponents
of $D$, $\tau^r$, and $\Gamma$, allowing any strictly larger
power for the final bound. Applied to the actual families,
\[
 \begin{array}{c|c|c}
 \text{family}&\text{multiplier and coefficient budget}&
                 \text{upper exponent}\\ \hline
 \text{cubic}&2h\,\tau^3\cdot C&-1/2\\
 \text{quartic}&2h\,\tau^4\cdot C&-1\\
 \text{quadratic loop}&\tau^2\cdot2P(h)&-1
 \end{array}
\]
Here $C$ is the incoming fixed path bound, $P(h)$ is the monitor,
and the channel weight has absolute value at most one. The
multiplier is the original $\beta\le2h$, wherever it occurs.
The fixed number of signed orientations and plaquette axes is
summed, not discarded.

For the tail of a squared path row, retain the ordered pair of
path labels $(i,k)$. With coefficients $a_i,a_k$ and lengths
$l_i,l_k$, its price is
\[
 b_{ik}=2|a_i a_k|\,
       1_{\{l_i d>R\ \mathrm{or}\ l_k d>R\}},
\]
where $d=1$ for the fine row and $d=L$ for the coarse row.
This bounds the actual half inner product because each unmasked
transported curvature has norm at most two; the kept-row
indicator can only decrease its magnitude. If
$\sum_i|a_i|e^{\sigma l_i}\le C$, then
\begin{equation}\label{eq:inherited-tail-budget}
 \sum_{i,k}b_{ik}e^{(\sigma/2)(l_i+l_k)}
 \le2e^{-\sigma R/(2d)}C^2.
\end{equation}
Indeed every retained pair has $l_i+l_k\ge R/d$; extract
$e^{-\sigma R/(2d)}$ and use the product of the two weighted
row sums. The genuine read support includes both full paths and
the kept-row test halo. After projection that extra halo is paid
by replacing $\ell$ by $\ell+2$. The site bound with its length
reserve is valid if $8\varepsilon<\sigma$, since its available
input exponent is now $\sigma/2$. As $R=\lfloor B/100000\rfloor$,
the right side of \eqref{eq:inherited-tail-budget} has every
negative upper exponent, even after multiplication by $2h$,
fixed powers of the true length and the site-conversion cost.
For the coarse-level tail the parent map is the identity; its
population factor is one, and the exponential still has fixed
$d=L$. The argument thus covers both tails without confusing
their spatial scales.

\subsection{Inherited jets in the compensated variation}
The incoming regular activity budget has upper exponent $-11/10$;
the incoming forced activity budget has every negative upper
exponent. The former bounds the original masked occurrence on
its healthy domain, extended by zero, and the latter bounds the
forced occurrence everywhere. Apply the anchor-to-site estimate
to their original supports and masks. This preserves those
exponents and all label multiplicities.

During the compensated variation the underlying fine field $U_*$
is held fixed. Hence the regular, cubic, quartic, loop, fine-tail
and forced values are constant in the coarse parameter. Their
protective factors vary. On a coarse support $Y$, each protective
bond cutoff has amplitude one and derivative scale of upper
exponent zero. Covering $Y$ by its $B$-cover, its full bond read
set contains at most
$b_B b_1\rho_L c_B(Y)$ bonds. The product rule therefore gives
the protector derivative scale $v_0(x)c_B(Y)$, with $v_0$ of
upper exponent zero. Multiplying a constant value by this
protector preserves its scalar price.

The coarse tail is the exception: its original coarse field varies.
For a path of length $l$, the unitary product argument gives
its transported curvature jets amplitude $|a_i|R_G^2$ and
derivative scale $2R_Gvl$. A product of two paths thus has scale
$2R_Gv(l_i+l_k)$, and its half inner product has amplitude at
most $|a_i a_k|R_G^4$. Multiplying by the original $\beta$
and channel coefficient is paid by
$2hR_G^4 b_{ik}$. The same tail estimate prices these labels
with every negative upper exponent, retaining the total length
in the derivative scale. There is no moving coarse kept-row
indicator in this term: its unmasked coarse counterpart is the
smooth path-pair expression with the fixed long-pair selector.

Choose $-1/2<t<0$. By the strict exponent slack, the six
inherited logarithmic families together have a rooted price
at most $h^t$ for sufficiently large $x$, with the support and
$2\varepsilon$ extra-cost reserves just proved. For the coarse
tail add its differentiated price to the scalar one; both are
nonnegative and negligible, so this addition is paid by choosing
an intermediate exponent between $-1/2$ and $t$ first.
The forced family has the same bound. Their common derivative
scale can be taken as
\begin{equation}\label{eq:inherited-jet-price}
 v_1(x)\{c_B(Y_i)+d_i\},\qquad
 v_1\text{ of upper exponent zero},
\end{equation}
where $d_i$ is the old support cover or the true path length,
as appropriate. The bounds are uniform over the held fine
configuration and preserve its original labels and parameters.
